\documentclass[11pt,a4paper]{article}
\usepackage{amsmath,amssymb,amsthm}
\usepackage[margin=2.5cm]{geometry}
\usepackage{hyperref}

\newtheorem{definition}{Definition}[section]
\newtheorem{theorem}[definition]{Theorem}
\newtheorem{proposition}[definition]{Proposition}
\newtheorem{lemma}[definition]{Lemma}
\newtheorem{corollary}[definition]{Corollary}
\newtheorem{conjecture}[definition]{Conjecture}
\newtheorem{remark}[definition]{Remark}

\title{Hyperseed Formalization:\\
  Three Things the World Doesn't Understand About AGI}
\author{ProtomegaTron (automated formalization)}
\date{2026-09-18}

\begin{document}
\maketitle

\begin{abstract}
We formalize the intellectual structure of Ben Goertzel's Substack article
``Three Things the World Doesn't Understand About AGI'' (2026-03-26) within
the Hyperseed ontology. Key mappings: the default LLM-scaling path as a
single-fiber assumption, alternative approaches as fiber diversity, Hyperon
as a heterogeneous fiber bundle, centralized vs.\ decentralized AGI as a
fiber topology choice, collective consciousness as fiber topology emergence,
organizational architecture affecting AGI quality, and the LLM/reasoning
split as base/fiber decomposition.
\end{abstract}

\tableofcontents

%% =================================================================
\section{Single-fiber assumption vs.\ fiber diversity}
\label{sec:paths}
%% =================================================================

\begin{theorem}[Thing 1: The default path --- claims 1--7, 19--20]
\label{thm:paths}
The ``just scale LLMs'' approach treats one fiber type (transformer
neural networks trained by backpropagation) as the only viable
cognitive architecture for AGI. This is a \emph{single-fiber assumption}:
it ignores the diversity of possible cognitive fiber types.

LLMs are fundamentally next-token predictors: they lack the ability
to truly reason about novel situations, reflect on their own thinking,
set and pursue goals, or build up understanding the way a developing
mind does. These are capacities that require different fiber types
--- different computational architectures with different structural
properties.

In Hyperseed terms, multiple fiber types exist for cognition:
\begin{itemize}
\item \textbf{Neural fiber} (LLMs, predictive coding): broad pattern
  recognition, knowledge retrieval, fluent generation. Strong at
  base-level operations (fast pattern matching over large domains).
\item \textbf{Symbolic fiber} (logic-based reasoning, MeTTa): precise
  deduction, knowledge manipulation, formal verification. Strong at
  fiber-level operations (structured reasoning about relationships).
\item \textbf{Evolutionary fiber} (evolutionary learning, genetic
  programming): open-ended exploration, creative search, adaptation.
  Strong at fiber-space exploration (finding novel fiber configurations).
\item \textbf{Probabilistic fiber} (probabilistic programming, PLN):
  uncertainty quantification, evidence accumulation, belief revision.
  Strong at fiber-level inference (reasoning under uncertainty).
\end{itemize}

Each fiber type has structural properties that make it suitable for
different aspects of general intelligence. No single fiber type is
sufficient --- AGI requires \emph{fiber diversity}: multiple fiber
types working together, each contributing its structural strengths.

Gary Marcus is closer to right than the ``just scale LLMs'' crowd,
but he overstates the case. LLMs aren't useless --- they're a
valuable fiber type that provides broad knowledge and fluent generation.
They're just not \emph{the only} fiber type needed, and not the most
important one for the capabilities that distinguish AGI from narrow AI.
\end{theorem}

%% =================================================================
\section{Heterogeneous fiber bundle}
\label{sec:hyperon}
%% =================================================================

\begin{proposition}[Hyperon architecture --- claims 4--6, 21, 25]
\label{prop:hyperon}
Hyperon integrates multiple AI paradigms --- neural, symbolic,
evolutionary, probabilistic --- in a common framework with MeTTa as
the connecting language.

In Hyperseed terms, Hyperon is a \emph{heterogeneous fiber bundle}:
multiple fiber types (neural, symbolic, evolutionary, probabilistic)
bundled together over a common base, connected through MeTTa as the
fiber-connecting infrastructure:
\[
  E_{\text{Hyperon}} = \bigcup_i F_i \quad \text{where } F_i \in
  \{\text{neural}, \text{symbolic}, \text{evolutionary},
  \text{probabilistic}\}
\]

MeTTa serves as the \emph{connection} on this bundle: it defines how
different fiber types communicate, how information flows between them,
how the results of one fiber type become inputs to another. The
connection is what makes the bundle more than a collection of
independent systems.

The LLM/reasoning split illustrates the base/fiber decomposition:
\begin{itemize}
\item \textbf{LLMs as base:} broad knowledge, fluent generation,
  pattern recognition --- the base over which reasoning operates.
  Wide but shallow.
\item \textbf{Reasoning (MeTTa, PLN) as fiber:} structured inference,
  logical deduction, belief revision --- the fiber that provides depth.
  Deep but narrow.
\end{itemize}

LLM-only approaches have base without fiber: broad knowledge without
deep reasoning. Hyperon provides both: the base for breadth (LLMs as
knowledge resource) and the fiber for depth (MeTTa/PLN for structured
reasoning). AGI requires the combination.
\end{proposition}

%% =================================================================
\section{Fiber topology choice: centralized vs.\ decentralized}
\label{sec:topology}
%% =================================================================

\begin{theorem}[Thing 2: Open-source matters --- claims 8--12, 22]
\label{thm:topology}
Whether AGI is built by centralized big-tech companies or by open-source
decentralized networks has enormous implications. The case for
decentralized AGI is structural, not just ideological: different
organizational architectures produce different failure modes, different
power dynamics, and different innovation patterns.

In Hyperseed terms, this is a \emph{fiber topology choice}:
\begin{itemize}
\item \textbf{Centralized AGI} = single controlled fiber with rigid
  boundary. One organization controls the fiber (the AGI system), decides
  who has access, determines how it's used. Location Zero organization:
  ego-driven, competitive, self-preserving.
\item \textbf{Decentralized AGI} = distributed fiber bundle with
  fork/merge topology. Multiple independent participants contribute to
  a shared fiber bundle. No single controller. Location 1+ organization:
  open, forkable, transparent, voluntary.
\end{itemize}

The structural differences:
\begin{center}
\begin{tabular}{ll}
\textbf{Centralized} & \textbf{Decentralized} \\
\hline
Single point of failure & No single point of failure \\
Power concentrated & Power distributed \\
Access controlled & Access open \\
Safety = small group's judgment & Safety = collective responsibility \\
Innovation = internal & Innovation = ecosystem-wide \\
Capture-prone & Capture-resistant \\
\end{tabular}
\end{center}

The ASI Alliance (SingularityNET, Fetch.ai, Ocean Protocol) is building
decentralized AGI infrastructure with blockchain-based coordination ---
a distributed fiber bundle with cryptographic coordination protocols
as the connection.
\end{theorem}

%% =================================================================
\section{Fiber topology emergence: collective consciousness}
\label{sec:consciousness}
%% =================================================================

\begin{definition}[Thing 3: Collective consciousness --- claims 13--18, 23--24]
\label{def:consciousness}
Jeffery Martin's consciousness locations can be applied to organizations.
Open-source organizations exhibit Location 1+ collective consciousness
not because their participants are individually enlightened but because
the organizational architecture produces emergent higher-consciousness
collective behavior.

In Hyperseed terms, collective consciousness is a \emph{fiber topology
emergence}: it depends on how fibers connect, how boundaries work,
how information flows --- the topology of the organizational fiber
bundle, not the content of individual fibers:
\[
  \text{Collective consciousness} = g(T_{\text{bundle}}) \neq
  \sum_i g(f_i)
\]

This has direct implications for AGI development: if organizational
architecture determines collective consciousness level, and collective
consciousness affects the quality and character of work, then the
architecture of AGI development organizations affects the AGI they
produce:
\[
  T_{\text{org}} \implies \text{collective consciousness} \implies
  \text{AGI character}
\]

A Location 1+ organization (open-source, decentralized, transparent)
may produce fundamentally different AGI than a Location Zero
organization (centralized, competitive, secretive). The AGI inherits
properties of the organizational fiber that produced it --- not
in a mystical sense but structurally: the development process shapes
the system through the decisions made, the priorities set, and the
values embedded in the architecture.

This connects Things 1, 2, and 3: the path to AGI (fiber diversity),
the organizational structure (fiber topology), and the collective
consciousness of the development community (fiber topology emergence)
are all aspects of the same underlying question --- what fiber bundle
will produce beneficial AGI?
\end{definition}

%% =================================================================
\section{Cross-references}
%% =================================================================

\begin{remark}[Connections to other formalizations]
\begin{itemize}
\item \textbf{The Architecture of Collective Non-Self} (2026-03-26):
  companion piece. Develops the ``Thing 3'' organizational consciousness
  analysis in detail.
\item \textbf{Open-Sourcing Your Self} (2026-03-27): extension. Takes
  the organizational consciousness analysis and applies it to mind
  uploading.
\item \textbf{Seeding RSI} (2026-07-28): Hyperon architecture in detail.
  The heterogeneous fiber bundle architecture applied to recursive
  self-improvement.
\item \textbf{OpenBGI} (2026-05-07): practical implementation. OpenBGI
  as a concrete instance of decentralized AGI governance architecture.
\item \textbf{The OpenAI/Microsoft Shakeup} (2026-05-01): negative
  example. Centralized AGI development (Location Zero organization)
  failing due to capital gravity --- exactly the risk identified here.
\item \textbf{Beyond Human Comparison} (2026-03-19): Four-Factor Model.
  What AGI actually is --- complementary to the organizational question
  of who builds it and how.
\end{itemize}
\end{remark}

\begin{conjecture}[Fiber bundle determines AGI character]
Conjecture: the character of an AGI system (its values, priorities,
behavioral tendencies) is determined more by the fiber bundle of the
organization that produced it than by the explicit alignment training
it received.

If this conjecture holds, alignment is primarily an organizational
architecture problem (designing the right fiber bundle for AGI
development) rather than primarily a technical problem (designing the
right training objective). The most important alignment decision is
not ``what objective function?'' but ``what organizational fiber
topology produces a development process that naturally produces
beneficial AGI?''

This doesn't eliminate the technical alignment problem --- it
reframes it as downstream of the organizational problem. The right
organizational fiber topology makes the technical alignment problem
easier because the development process naturally prioritizes the
right things.
\end{conjecture}

\end{document}
